\documentclass[10pt,a4paper]{amsart}
\usepackage[margin=19mm]{geometry}
\usepackage{amsmath,amssymb,mathtools}
\usepackage[scr=rsfs,cal=euler]{mathalfa}
\usepackage{xfrac}
\usepackage[unicode,hidelinks,bookmarks=false]{hyperref}
\newcommand{\ie}{i.e.\ }
\newcommand{\cf}{cf.\ }
\newcommand{\resp}{respectively\ }
\newcommand{\Subsection}{\subsection}
\providecommand{\nobreakdash}{\nobreak\hbox{-}}
\setlength{\emergencystretch}{2em}
\multlinegap0pt
\usepackage{amssymb}
\usepackage[shortlabels]{enumitem}
\usepackage{csquotes}
\usepackage{mathtools}
\newtheorem{lemma}{Lemma}
 \newcommand{\R}{\mathbb R}
 \DeclareMathOperator{\Real}{Re}
 \newcommand{\dd}{\mathrm d}
\title{On the propagation of mountain waves: Linear theory}
\author{Adrian Constantin, J\"org Weber}
\date{Source: arXiv:2509.26125v2 (CC BY 4.0)}
\begin{document}
\maketitle

\noindent\emph{Source and licence.} On the propagation of mountain waves: Linear theory. Version arXiv:2509.26125v2 (CC BY 4.0), distributed by arXiv under the Creative Commons Attribution 4.0 International licence, http://creativecommons.org/licenses/by/4.0/, as recorded in the arXiv licence metadata for this exact version and shown on the abstract page https://arxiv.org/abs/2509.26125v2. Full licence notice: \texttt{licenses/arxiv-2509.26125-CC-BY-4.0.txt}.\par\smallskip

\noindent\textit{Editorial scope.} Counted unit: the article's lemma lem:K\_est (source0.tex lines 711--715). The article's complete original proof of it is delivered with it (source0.tex lines 716--755, one \texttt{proof} environment of 40 lines). No mathematics is written, repaired or re-derived: the statement and the proof are the article's own lines, in the article's own notation, and no retained line is substituted or re-flowed.  Retained uncounted as the source-local context the proof needs: lines 421--424; lines 555--558; lines 696--703.  Nothing was omitted from any retained span, so the package carries no omission record.  No retained line contains a citation, so the wrapper carries no bibliography and no bibliography entry.  No retained line contains a figure environment, an image-inclusion command or drawing code, so no image is delivered and the package declares no asset dependency.  Editorial formatting only, in addition: an AMS \texttt{amsart} preamble that restates the article's own declarations and macros used by the retained lines, namely the environments \texttt{lemma} on separate counters, together with the standard packages those lines need.  The retained environments print in this document as Lemma 1 in order of appearance, whereas the article prints the same items with the numbering of its own full document.  The complete native source of the article as distributed in the arXiv e-print for this exact version is bundled unchanged as \texttt{sources/186060/source0.tex}.
\begin{align}
	\text{there exists }F_0 > 0 \text{ such that }& F-F_0, F' \in L^1([0,\infty))\nonumber\\
	\text{and }&\int_0^\infty z|F(z)-F_0|\,\dd z<\infty.\tag{A}\label{ass:A}
\end{align}

\begin{align}
	K_0(x,z) &= \int_0^\infty \frac{e^{-\sqrt\lambda|x|}}{2\sqrt\lambda}\frac{\sin(\sqrt\lambda z)}{\sqrt\lambda}\frac{\sqrt\lambda}{\pi}\,\dd\lambda = \frac{1}{\pi} \int_0^\infty e^{-\mu|x|}\sin(\mu z)\,\dd \mu \nonumber\\
	&= \frac{z}{\pi(x^2+z^2)}.\label{eq:K0}
\end{align}

	\begin{align}
		\sigma &\in C((0,\infty)),\label{eq:sigma_continuous}\\
		\sigma(\lambda) &= \mathcal O\left(\frac{1}{\sqrt\lambda}\right),\quad \lambda\to 0,\label{eq:sigma_asymptotics_zero}\\
		\sigma(\lambda) &= \frac{\sqrt\lambda}{\pi}+\mathcal O\left(\frac{1}{\sqrt\lambda}\right),\quad \lambda\to\infty,\label{eq:sigma_asymptotics_infinity}\\
		|v(z,\lambda)| &\le \frac{c}{\sqrt\lambda},\label{eq:est_v}\\
		v(z,\lambda)-\frac{\sin(\sqrt\lambda z)}{\sqrt\lambda} &= -\frac{\cos(\sqrt\lambda z)}{2\lambda}\int_0^z q(z')\,\dd z'\, + \mathcal O(\lambda^{-3/2}),\quad \lambda\to\infty,\label{eq:asymp_v-v0}\\
		|v(z,\lambda)\sigma(\lambda)| &\le \begin{cases}c,&\lambda\ge1,\\\frac{c}{\sqrt\lambda},&0<\lambda\le1.\end{cases}\label{eq:est_v*sigma}
	\end{align}

\begin{lemma}\label{lem:K_est}
	Under Assumption \eqref{ass:A}, the singular part of $K$ is precisely $K_0$ in \eqref{eq:K0}, in the sense that $K-K_0$ is continuous on $\overline{\R^2_+}$ and
	\[|K(x,z)-K_0(x,z)| \le c\]
	with $c>0$ independent of $x$ and $z$.
\end{lemma}

\begin{proof}
	The deviation of the evanescent from the unperturbed piece can be written as follows:
	\begin{align*}
		&K^{\mathrm{e}}(x,z) - K_0(x,z) \\
		&= \int_{F_0}^\infty\frac{e^{-\sqrt{\lambda-F_0}|x|}}{2\sqrt{\lambda-F_0}}v(z,\lambda)\sigma(\lambda)\,\dd\lambda - \int_0^\infty \frac{e^{-\sqrt\lambda|x|}}{2\sqrt\lambda}\frac{\sin(\sqrt\lambda z)}{\sqrt\lambda}\frac{\sqrt\lambda}{\pi}\,\dd\lambda\\
		&= \int_{F_0}^\infty\left(\frac{e^{-\sqrt{\lambda-F_0}|x|}}{2\sqrt{\lambda-F_0}}-\frac{e^{-\sqrt\lambda|x|}}{2\sqrt\lambda}\right)v(z,\lambda)\sigma(\lambda)\,\dd\lambda\\
		&\quad\;+\int_{F_0}^\infty \frac{e^{-\sqrt\lambda|x|}}{2\sqrt\lambda}v(z,\lambda)\left(\sigma(\lambda)-\frac{\sqrt\lambda}{\pi}\right)\,\dd\lambda\\
		&\quad\;+\int_{F_0}^\infty \frac{e^{-\sqrt\lambda|x|}}{2\sqrt\lambda}\left(v(z,\lambda)-\frac{\sin(\sqrt\lambda z)}{\sqrt\lambda}\right)\frac{\sqrt\lambda}{\pi}\,\dd\lambda\\
		&\quad\;-\int_0^{F_0} \frac{e^{-\sqrt\lambda|x|}}{2\sqrt\lambda}\frac{\sin(\sqrt\lambda z)}{\sqrt\lambda}\frac{\sqrt\lambda}{\pi}\,\dd\lambda\\
		&\eqqcolon I_1+I_2+I_3+I_4.
	\end{align*}
	First, since $\dd(e^{-y}/y)/\dd y=-e^{-y}(y+1)/y^2$, we have
	\begin{align*}
		\left|\frac{e^{-\sqrt{\lambda-F_0}|x|}}{2\sqrt{\lambda-F_0}}-\frac{e^{-\sqrt\lambda|x|}}{2\sqrt\lambda}\right| &\le \frac{e^{-\sqrt{\lambda-F_0}|x|}(\sqrt{\lambda-F_0}|x|+1)}{2(\lambda-F_0)}\, (\sqrt\lambda-\sqrt{\lambda-F_0}) \\
		&\le c\lambda^{-3/2}
	\end{align*}
	for $\lambda\ge F_0+1$ independent of $x$. Thus and by \eqref{eq:est_v*sigma},
	\[|I_1| \le \int_{F_0}^{F_0+1} \frac{c}{\sqrt{\lambda-F_0}}\,\dd\lambda + c\int_{F_0+1}^\infty \lambda^{-3/2} \,\dd\lambda < \infty.\]
	Second, by \eqref{eq:sigma_asymptotics_infinity} and \eqref{eq:est_v},
	\[|I_2| \le c \int_{F_0}^\infty \lambda^{-3/2}\,\dd\lambda < \infty.\]
	Third, by \eqref{eq:est_v} and \eqref{eq:asymp_v-v0},
	\[\left|I_3 + \frac{1}{4\pi} \int_0^z q(z')\,\dd z' \int_{F_0}^\infty \frac{e^{-\sqrt\lambda|x|} \cos(\sqrt\lambda z)}{\lambda} \,\dd\lambda\right| \le c \int_{F_0}^\infty \lambda^{-3/2}\,\dd\lambda < \infty.\]
	Here, we realize that
	\[\int_1^\infty \frac{e^{-\sqrt\lambda|x|} \cos(\sqrt\lambda z)}{\lambda} \,\dd\lambda = 2\Real \int_1^\infty \frac{e^{\mu(-|x|+iz)}}{\mu} \,\dd\mu.\]
	It is well-known that this exponential integral has a logarithmic singularity near $(x,z)=(0,0)$ and, away from $(0,0)$, is bounded. Therefore and by $q\in L^1([0,\infty))\cap L^\infty([0,\infty))$,
	\[|I_3| \le c + c\left|\int_0^z q(z')\,\dd z' \, \Real \int_1^\infty \frac{e^{\mu(-|x|+iz)}}{\mu} \,\dd\mu\right| \le c.\]
	Fourth, clearly
	\[|I_4| \le c.\]
	Summing up, we have proved that
	\[|K^{\mathrm{e}}(x,z) - K_0(x,z)| \le c,\]
	where $c$ is independent of $x$ and $z$.
	
	The radiated and trapped pieces of $K$ are fortunately much easier to handle. Indeed, by \eqref{eq:est_v*sigma} we have
	\[|K^{\mathrm{r}}(x,z)|\le c\int_0^{F_0} \frac{\dd\lambda}{\sqrt{F_0-\lambda} \sqrt\lambda} <\infty\]
	and, by just being a finite sum,
	\[|K^{\mathrm{t}}(x,z)|\le c\]
	in view of \eqref{eq:est_v}.
	
	The continuity of $K-K_0$ follows directly from the above estimates.
\end{proof}

\end{document}
